% ---------------------------------------------------------------------
% Preâmbulo: pacotes, geometria, cores, macros e ambientes
% ---------------------------------------------------------------------
\usepackage[brazil,english]{babel}
\usepackage[utf8]{inputenc}
\usepackage[T1]{fontenc}
\usepackage{amsmath, amssymb, amsfonts}
\usepackage{graphicx}
\usepackage{geometry}
\usepackage{enumitem}
\usepackage{multicol}
\usepackage{xcolor}
\usepackage{colortbl}
\usepackage{setspace}
\usepackage{array}
\usepackage{longtable}
\usepackage{booktabs}
\usepackage{multirow}
\usepackage{tikz}
\usetikzlibrary{arrows.meta, positioning, fit, backgrounds}
\usepackage{hyperref}

\geometry{top=2.5cm, bottom=2.5cm, left=2.5cm, right=2.5cm}

% ---------- Cores ----------
\definecolor{pucazul}{HTML}{0B5A85}
\definecolor{cinzaclaro}{HTML}{EEF2F6}
\definecolor{destaque}{HTML}{DCEAF5}

\hypersetup{
    colorlinks=true,
    linkcolor=pucazul,
    urlcolor=pucazul,
    pdftitle={Canais, estratégia e alinhamento com a proposta de valor -- PetCuida},
    pdfauthor={Grupo 4 -- Tópicos III},
}

% ---------- Espaçamento ----------
\setlength{\parindent}{1.25cm}
\setlength{\parskip}{4pt}
\setlength{\emergencystretch}{2em}
\renewcommand{\arraystretch}{1.25}
\setlist{itemsep=2pt, topsep=4pt}

% ---------- Dados do trabalho (edite aqui, reflete na capa e folha de rosto) ----------
\newcommand{\projeto}{PetCuida}
\newcommand{\tituloatividade}{Canais, estratégia e alinhamento com a proposta de valor}
\newcommand{\tipoatividade}{(ATIVIDADE EM GRUPO)}
\newcommand{\anoentrega}{2026}
\newcommand{\orientador}{João Carlos Oliveira Caetano}
\newcommand{\tutor}{Paulo Mattos}

% ---------- Ambientes de elementos pré-textuais ----------
\newenvironment{capa}{\thispagestyle{empty}}{\clearpage}
\newenvironment{folhaderosto}{\thispagestyle{empty}}{\clearpage}

% ---------- Macros de conteúdo ----------
% Marca de canal priorizado
\newcommand{\prio}{\textcolor{pucazul}{$\bigstar$}}
% Nome de canal em negrito
\newcommand{\canal}[1]{\textbf{#1}}

% Colunas de tabelas longas (p{} com alinhamento à esquerda)
\newcolumntype{L}[1]{>{\raggedright\arraybackslash}p{#1}}
\newcolumntype{C}[1]{>{\centering\arraybackslash}p{#1}}

% Cabeçalho de tabela com fundo colorido
\newcommand{\cab}[1]{\textbf{\color{white}#1}}
\newcommand{\cabfundo}{\rowcolor{pucazul}}

% Justificativa de um canal: \justcanal{Nome}{alinhamento}{adequação}{viabilidade}
\newcommand{\justcanal}[4]{%
  \subsubsection{#1}
  \begin{description}[font=\bfseries\itshape, leftmargin=1.2em, style=unboxed, itemsep=3pt]
    \item[Alinhamento com a proposta de valor.] #2
    \item[Adequação ao cliente-alvo.] #3
    \item[Viabilidade para a startup.] #4
  \end{description}
}
